\abstrakinggris{Nasi Gerilya Restaurant is a Padang food business in Medan that uses GrabFood as one of its main sales channels, yet the situation at the research site shows unstable sales, a declining proportion of repeat purchases, customer complaints, and competitive pressure from similar merchants on the same platform. The purpose of this study is to identify the internal and external factors of Nasi Gerilya Restaurant on the GrabFood platform and to formulate marketing strategies based on SWOT analysis to support sales improvement. The research was conducted at Nasi Gerilya Restaurant in Medan, focusing on its marketing activities through GrabFood. Data were collected through interviews with the owner, employees, customers, and comparative informants, GrabFood platform observation, customer questionnaires, business documentation, customer reviews, and literature study. The results show that an IFAS score of 2.50 and an EFAS score of 2.51 place Nasi Gerilya in a moderate condition and in Quadrant I, namely an aggressive or SO strategy position, because internal strengths and external opportunities are more dominant than weaknesses and threats. The main strengths are strong food taste, large portions, digital reputation, an order-checking flow, and complaint response, while the main weaknesses are packing accuracy, cross-channel queues, stock updates, price perception, and unclear package descriptions. The alternative strategies formulated include SO, WO, ST, and WT strategies, with priorities on improving the queue system and packing checklist, strengthening segment-based hero menu packages, improving the digital storefront and digital menu, synchronizing promotions with stock, capacity, cost of goods sold, and margins, and utilizing digital stamps, customer databases, and Flowie to encourage repeat purchases.}

\katakunciinggris{marketing strategy, GrabFood, SWOT, culinary MSME, sales}
